\documentclass[11pt]{amsart}
\usepackage[T1]{fontenc}
\usepackage[utf8]{inputenc}
\usepackage{lmodern}
\usepackage{amsmath}
\usepackage{amssymb}
\usepackage{amsfonts}
\usepackage{amsthm}
\usepackage{geometry}
\geometry{margin=1in}
\numberwithin{equation}{section}
\setcounter{tocdepth}{3}
\setcounter{secnumdepth}{3}
\theoremstyle{plain}
\newtheorem{theorem}{Theorem}[section]
\newtheorem{proposition}[theorem]{Proposition}
\newtheorem{lemma}[theorem]{Lemma}
\newtheorem{corollary}[theorem]{Corollary}
\newtheorem{conjecture}[theorem]{Conjecture}
\theoremstyle{definition}
\newtheorem{definition}[theorem]{Definition}
\newtheorem{experiment}[theorem]{Experiment}
\newtheorem{heuristic}[theorem]{Heuristic}
\newtheorem{question}[theorem]{Question}
\theoremstyle{remark}
\newtheorem{remark}[theorem]{Remark}
\newtheorem{warning}[theorem]{Warning}
\usepackage[hidelinks]{hyperref}

\title[Moving Rigid Leslie Models]{Moving Rigid Leslie Models on Consecutive Logarithmic Windows:\\Critical Spectral Geometry, Overlap Rigidity, Scalar Innovations,\\and a Causal Forecasting Test}
\author{}
\date{}

\begin{document}
\maketitle
\tableofcontents
\bigskip

\begin{remark}
\label{rem:status}
This document is a research notebook.  The algebraic identities and local differential constraints are proved under explicit regularity and simplicity hypotheses.  The proposed connection with consecutive windows of the Mertens function is an experimental programme, not a proved arithmetic theorem.  No statement below proves a new estimate for the Mertens function or the Riemann hypothesis.
\end{remark}

\section{Purpose and change of viewpoint}
\label{sec:purpose}

The earlier static experiments asked whether one finite collection of fitted frequencies could be embedded in one nonnegative Leslie matrix.  Allowing auxiliary eigenvalues made that static problem feasible in all forty reported configurations, but the result did not distinguish a coherent dynamical law from independent fitting.

The present formulation replaces independent static realizations by a sequence of rigid Leslie matrices attached to consecutive logarithmic windows.  The object of study is therefore a non-autonomous product, or linear cocycle, rather than one fixed matrix.  Linear cocycles and their long products are standard objects in dynamical-systems theory \cite{DuarteKlein2016}.

\begin{warning}
\label{warn:no_automatic_novelty}
The use of a time-dependent matrix product is not itself a new result.  The potentially discriminating content is the simultaneous enforcement of rigidity, nonnegativity, continuity through the feasible set, and prediction of a future window without recalibration.
\end{warning}

\section{Rigid moving Leslie matrices}
\label{sec:rigid_moving_leslie}

\begin{definition}
\label{def:leslie_matrix}
For $m\geq2$, a companion Leslie matrix is
\begin{equation}
L(a)=
\begin{pmatrix}
a_1&a_2&\cdots&a_m\\
1&0&\cdots&0\\
0&1&\ddots&\vdots\\
\vdots&&\ddots&0
\end{pmatrix},
\qquad a_\ell\geq0.
\label{eq:leslie_matrix}
\end{equation}
Its characteristic polynomial is
\begin{equation}
P(z;a)=z^m-\sum_{\ell=1}^{m}a_\ell z^{m-\ell}.
\label{eq:characteristic_polynomial}
\end{equation}
\end{definition}

Fix $K\geq1$ and set
\begin{equation}
m=2K+1.
\label{eq:rigid_dimension}
\end{equation}
For $b>1$ and phases $\Omega=(\omega_1,\ldots,\omega_K)$, prescribe the roots
\begin{equation}
b,
\qquad
\sqrt b\,e^{\pm i\omega_1},\ldots,
\sqrt b\,e^{\pm i\omega_K}.
\label{eq:prescribed_roots}
\end{equation}
The unique monic polynomial with these roots is
\begin{equation}
P_{b,\Omega}(z)
=(z-b)\prod_{k=1}^{K}
\left(z^2-2\sqrt b\cos(\omega_k)z+b\right).
\label{eq:rigid_polynomial}
\end{equation}
Writing
\begin{equation}
P_{b,\Omega}(z)
=z^m-\sum_{\ell=1}^{m}a_\ell(b,\Omega)z^{m-\ell}
\label{eq:rigid_coefficients}
\end{equation}
defines the coefficient map $A:(b,\Omega)\mapsto a(b,\Omega)$.

\begin{proposition}
\label{prop:uniqueness}
For fixed $(b,\Omega)$, the rigid Leslie coefficients in \eqref{eq:rigid_coefficients} are unique.  A rigid nonnegative Leslie realization exists if and only if
\begin{equation}
a_\ell(b,\Omega)\geq0,
\qquad 1\leq\ell\leq m.
\label{eq:coefficient_nonnegativity}
\end{equation}
\end{proposition}

\begin{proof}
The roots in \eqref{eq:prescribed_roots} determine the monic polynomial \eqref{eq:rigid_polynomial}.  Comparison with \eqref{eq:characteristic_polynomial} determines every coefficient.  The companion matrix is entrywise nonnegative exactly when its first-row coefficients satisfy \eqref{eq:coefficient_nonnegativity}.
\end{proof}

\begin{definition}
\label{def:admissible_phase_set}
The rigid admissible phase set is
\begin{equation}
\mathcal A_{b,K}
=
\left\{
\Omega\in\mathbb T^K:
 a_\ell(b,\Omega)\geq0
 \text{ for all }\ell
\right\}.
\label{eq:admissible_phase_set}
\end{equation}
\end{definition}

\begin{remark}
\label{rem:boundary}
The boundary of \eqref{eq:admissible_phase_set} is contained in the union of the coefficient zero sets
\begin{equation}
\partial\mathcal A_{b,K}
\subseteq
\bigcup_{\ell=1}^{m}
\left\{\Omega:a_\ell(b,\Omega)=0\right\}.
\label{eq:admissible_boundary}
\end{equation}
Equality need not hold without excluding redundant inequalities and singular intersections.
\end{remark}

\section{Windowed cocycle and accumulated phase}
\label{sec:windowed_cocycle}

Let $W_0,W_1,\ldots$ be consecutive windows in logarithmic scale.  Associate a rigid Leslie matrix
\begin{equation}
L_j=L(a_j),
\qquad
a_j=a(b_j,\Omega_j)\in\mathbb R_+^m.
\label{eq:window_leslie}
\end{equation}
The non-autonomous evolution is
\begin{equation}
v_{j+1}=L_jv_j,
\label{eq:nonautonomous_evolution}
\end{equation}
with cocycle product
\begin{equation}
\Phi(J,0)=L_{J-1}\cdots L_1L_0.
\label{eq:cocycle_product}
\end{equation}

\begin{remark}
\label{rem:no_forced_periodicity}
A moving sequence $L_j$ has no mandatory spectral periodicity.  This does not imply that every chosen trajectory is nonperiodic.  Periodicity may reappear if the coefficients or phases are selected periodically.
\end{remark}

If a locally matched branch has phase $\omega_{j,k}$ in window $j$, its formal accumulated phase is
\begin{equation}
\Theta_k(J)=\sum_{j=0}^{J-1}\omega_{j,k}.
\label{eq:accumulated_phase}
\end{equation}
The definition requires mode matching and phase unwrapping between adjacent windows.  Without those operations, the unordered roots of successive polynomials do not determine a unique trajectory $k\mapsto\omega_{j,k}$.

\section{Exact coefficient coordinates for the rigid critical spectrum}
\label{sec:eta_coordinates}

The general inverse spectral problem for Leslie matrices has an established characterization in the literature \cite{Benvenuti2019}.  The results in this section concern the special rigid spectrum in \eqref{eq:prescribed_roots} and its use as a moving-window coordinate system.

Put $r=\sqrt b$ and define
\begin{equation}
S_\Omega(w)
=\prod_{k=1}^{K}\left(w^2-2\cos(\omega_k)w+1\right)
=\sum_{j=0}^{2K}d_jw^{2K-j}.
\label{eq:unit_circle_factor}
\end{equation}
The polynomial is reciprocal, and hence
\begin{equation}
d_{2K-j}=d_j,
\qquad
d_0=d_{2K}=1.
\label{eq:palindromic_d}
\end{equation}
Moreover,
\begin{equation}
P_{b,\Omega}(z)
=(z-r^2)r^{2K}S_\Omega(z/r).
\label{eq:scaled_rigid_polynomial}
\end{equation}

\begin{proposition}
\label{prop:coefficient_formula}
The coefficients in \eqref{eq:rigid_coefficients} satisfy
\begin{align}
a_\ell
&=r^\ell\left(rd_{\ell-1}-d_\ell\right),
&&1\leq\ell\leq2K,
\label{eq:coefficient_formula_main}\\
a_m&=r^{m+1}=b^{K+1}.
\label{eq:last_coefficient_fixed}
\end{align}
For $1\leq j\leq K$, the paired coefficients are
\begin{align}
a_j
&=r^j\left(rd_{j-1}-d_j\right),
\label{eq:left_paired_coefficient}\\
a_{m-j}
&=r^{m-j}\left(rd_j-d_{j-1}\right).
\label{eq:right_paired_coefficient}
\end{align}
\end{proposition}

\begin{proof}
Expand \eqref{eq:scaled_rigid_polynomial} and compare the coefficient of $z^{m-\ell}$ with \eqref{eq:rigid_coefficients}.  Equation \eqref{eq:right_paired_coefficient} follows from \eqref{eq:palindromic_d}.
\end{proof}

\begin{corollary}
\label{cor:ratio_test}
A rigid nonnegative Leslie realization exists if and only if
\begin{equation}
d_j>0
\quad\text{and}\quad
\frac1r\leq\frac{d_j}{d_{j-1}}\leq r,
\qquad 1\leq j\leq K.
\label{eq:ratio_test}
\end{equation}
Strict inequalities are equivalent to strict positivity of all paired coefficients.
\end{corollary}

\begin{proof}
Equations \eqref{eq:left_paired_coefficient} and \eqref{eq:right_paired_coefficient} give the two inequalities.  If one of $d_{j-1}$ and $d_j$ vanished while both paired inequalities held, the other would also vanish.  Propagation back to $d_0=1$ yields a contradiction.  Thus feasibility forces positivity of all $d_j$ for $0\leq j\leq K$.
\end{proof}

\begin{definition}
\label{def:eta_coordinates}
On the region where $d_1,\ldots,d_K>0$, define
\begin{equation}
\eta_j(\Omega)=\log\frac{d_j}{d_{j-1}},
\qquad 1\leq j\leq K.
\label{eq:eta_definition}
\end{equation}
The rigid spectral strain is
\begin{equation}
\Sigma(\Omega)=\|\eta(\Omega)\|_\infty.
\label{eq:spectral_strain}
\end{equation}
\end{definition}

\begin{corollary}
\label{cor:box_criterion}
The positivity test becomes
\begin{equation}
\|\eta(\Omega)\|_\infty
\leq\frac12\log b.
\label{eq:eta_box}
\end{equation}
The minimum compatible Perron value is
\begin{align}
b_*(\Omega)
&=\exp\left(2\|\eta(\Omega)\|_\infty\right)
\label{eq:minimal_perron_eta}\\
&=\max_{1\leq j\leq K}
\left\{
\left(\frac{d_j}{d_{j-1}}\right)^2,
\left(\frac{d_{j-1}}{d_j}\right)^2
\right\}.
\label{eq:minimal_perron_ratios}
\end{align}
If some $d_j\leq0$, set $b_*(\Omega)=+\infty$.
\end{corollary}

\begin{remark}
\label{rem:minimal_perron_boundary}
If $b>b_*(\Omega)$, the rigid Leslie coefficients are strictly positive.  If $b=b_*(\Omega)<\infty$, at least one paired coefficient vanishes.  The map $b_*$ is generally nonsmooth where several coordinates attain the maximum in \eqref{eq:minimal_perron_eta}.
\end{remark}

\begin{proposition}
\label{prop:comb_zero_strain}
The rigid root-of-unity comb
\begin{equation}
\omega_k=\frac{2\pi k}{2K+1},
\qquad 1\leq k\leq K,
\label{eq:comb_phases}
\end{equation}
has $d_j=1$ for all $j$, and hence $\Sigma(\Omega)=0$.  Conversely, $\Sigma(\Omega)=0$ implies the same multiset of phases, up to ordering and conjugation.
\end{proposition}

\begin{proof}
For \eqref{eq:comb_phases},
\begin{equation}
S_\Omega(w)=\frac{w^{2K+1}-1}{w-1}
=1+w+\cdots+w^{2K}.
\label{eq:comb_polynomial}
\end{equation}
Conversely, $\eta=0$ and $d_0=1$ imply $d_0=\cdots=d_{2K}=1$.  The roots of \eqref{eq:comb_polynomial} are the nontrivial $(2K+1)$-st roots of unity.
\end{proof}

\section{Memory tails and the universal critical region}
\label{sec:memory_critical_region}

Normalize the coefficients by
\begin{equation}
x_\ell=\frac{a_\ell}{b^\ell},
\qquad
X_q=\sum_{\ell=q}^{m}x_\ell.
\label{eq:memory_tail_definition}
\end{equation}

\begin{proposition}
\label{prop:memory_tail_formula}
The tails telescope to
\begin{equation}
X_{j+1}=\frac{d_j}{r^j},
\qquad 0\leq j\leq2K.
\label{eq:memory_tail_formula}
\end{equation}
Consequently,
\begin{align}
X_{m-j}&=b^{j-K}X_{j+1},
&&0\leq j\leq K,
\label{eq:tail_palindromy}\\
x_m&=X_m=b^{-K}.
\label{eq:last_memory_mass}
\end{align}
Moreover, positivity is equivalent to
\begin{equation}
\frac{X_j}{b}\leq X_{j+1}\leq X_j,
\qquad 1\leq j\leq K.
\label{eq:tail_inequalities}
\end{equation}
\end{proposition}

\begin{proof}
Substitute \eqref{eq:coefficient_formula_main} into \eqref{eq:memory_tail_definition}; consecutive terms telescope.  Equations \eqref{eq:tail_palindromy}--\eqref{eq:last_memory_mass} follow from \eqref{eq:palindromic_d}, and \eqref{eq:tail_inequalities} is \eqref{eq:ratio_test} rewritten using \eqref{eq:memory_tail_formula}.
\end{proof}

The reciprocal polynomial admits a Chebyshev reduction.  Define
\begin{equation}
C_n(y)=2T_n(y/2),
\label{eq:scaled_chebyshev}
\end{equation}
where $T_n$ is the Chebyshev polynomial of the first kind, and set
\begin{equation}
H_\eta(y)
=C_K(y)+d_1C_{K-1}(y)+\cdots+d_{K-1}C_1(y)+d_K.
\label{eq:H_eta}
\end{equation}
Then
\begin{equation}
S_\Omega(w)=w^KH_\eta(w+w^{-1}).
\label{eq:chebyshev_reduction}
\end{equation}

\begin{definition}
\label{def:universal_region}
Let
\begin{equation}
\mathcal R_K
=
\left\{
\eta\in\mathbb R^K:
H_\eta\text{ has }K\text{ simple roots in }(-2,2)
\right\},
\label{eq:universal_region}
\end{equation}
where $d_j=\exp(\eta_1+\cdots+\eta_j)$.
\end{definition}

\begin{proposition}
\label{prop:separation_geometry_positivity}
In $\eta$-coordinates, the admissible rigid critical region is
\begin{equation}
\mathcal A^{\eta}_{b,K}
=
\mathcal R_K
\cap
\left[-\frac12\log b,\frac12\log b\right]^K.
\label{eq:critical_box_intersection}
\end{equation}
Thus $\mathcal R_K$ encodes the unit-circle root geometry, while the box encodes Leslie positivity.
\end{proposition}

\begin{remark}
\label{rem:two_boundary_mechanisms}
Within the positive-$d$ coordinates, a path can leave \eqref{eq:critical_box_intersection} either through a box face or through the boundary of $\mathcal R_K$.  The latter occurs when roots collide or reach an endpoint.  Since
\begin{equation}
H_\eta(2)=d_K+2\sum_{j=0}^{K-1}d_j>0,
\label{eq:H_at_plus_two}
\end{equation}
a root cannot reach $+2$.  Candidate spectral boundary equations are therefore
\begin{equation}
\operatorname{Disc}H_\eta=0
\qquad\text{or}\qquad
H_\eta(-2)=0.
\label{eq:critical_region_boundary}
\end{equation}
This describes the regular boundary mechanisms; intersections and higher degeneracies form singular strata.
\end{remark}

\section{Jacobian of the phase-to-strain map}
\label{sec:eta_jacobian}

Assume distinct phases $0<\omega_k<\pi$.  Let the phases be ordered to remove permutation ambiguity.

\begin{proposition}
\label{prop:d_jacobian}
The coefficient Jacobian satisfies
\begin{equation}
\left|
\det\frac{\partial(d_1,\ldots,d_K)}
{\partial(\omega_1,\ldots,\omega_K)}
\right|
=
2^{K(K+1)/2}
\prod_{k=1}^{K}\sin\omega_k
\prod_{i<j}|\cos\omega_i-\cos\omega_j|.
\label{eq:d_jacobian}
\end{equation}
Consequently,
\begin{equation}
\left|
\det\frac{\partial\eta}{\partial\Omega}
\right|
=
\frac{
2^{K(K+1)/2}
\displaystyle\prod_k\sin\omega_k
\displaystyle\prod_{i<j}|\cos\omega_i-\cos\omega_j|
}{\displaystyle\prod_{j=1}^{K}d_j}.
\label{eq:eta_jacobian}
\end{equation}
\end{proposition}

\begin{proof}
The first $K$ reciprocal coefficients determine the monic polynomial with roots $2\cos\omega_k$ in the Chebyshev basis.  The root-to-coefficient Jacobian is a Vandermonde product, the change $2\cos\omega_k$ contributes $2\sin\omega_k$ in absolute value, and the triangular changes between monomial, Chebyshev, and reciprocal coefficient coordinates give the stated power of $2$.  The map $d\mapsto\eta$ is triangular with determinant $(d_1\cdots d_K)^{-1}$.
\end{proof}

\begin{corollary}
\label{cor:discriminant_jacobian}
The same determinant can be written as
\begin{equation}
\left|
\det\frac{\partial\eta}{\partial\Omega}
\right|
=
\frac{
\sqrt{\left|
\operatorname{Disc}H_\eta\,
H_\eta(2)H_\eta(-2)
\right|}
}{d_1d_2\cdots d_K}.
\label{eq:eta_jacobian_discriminant}
\end{equation}
Hence the phase-to-strain map is locally invertible away from phase collisions and the endpoints $0$ and $\pi$.
\end{corollary}

\begin{remark}
\label{rem:reduced_transport_coordinates}
At regular points, a moving rigid Leslie model can be transported using $K$ variables $\eta$, plus $b$ when the Perron value is allowed to vary.  The full coefficient vector is reconstructed from \eqref{eq:left_paired_coefficient}, \eqref{eq:right_paired_coefficient}, and \eqref{eq:last_coefficient_fixed}.
\end{remark}

\section{Exact viability conditions in strain coordinates}
\label{sec:eta_viability}

For fixed $b$, the box criterion \eqref{eq:eta_box} gives the exact first-order viability conditions.  On an upper active face,
\begin{equation}
\eta_j=\frac12\log b
\quad\Longrightarrow\quad
\dot\eta_j\leq0,
\label{eq:upper_face_velocity}
\end{equation}
and on a lower active face,
\begin{equation}
\eta_j=-\frac12\log b
\quad\Longrightarrow\quad
\dot\eta_j\geq0.
\label{eq:lower_face_velocity}
\end{equation}
Since $\dot\eta=D_\Omega\eta\,\dot\Omega$, both cases are summarized by
\begin{equation}
\operatorname{sgn}(\eta_j)
\sum_{k=1}^{K}
\frac{\partial\eta_j}{\partial\omega_k}
\dot\omega_k
\leq0
\label{eq:phase_velocity_face_condition}
\end{equation}
on an active face.

If $b=b(t)$ varies and
\begin{equation}
\eta_j=s_j\frac12\log b,
\qquad s_j\in\{-1,1\},
\label{eq:moving_active_face}
\end{equation}
then viability requires
\begin{equation}
s_j\dot\eta_j
\leq\frac12\frac{\dot b}{b}.
\label{eq:moving_face_velocity}
\end{equation}
For the minimum-Perron choice $b=b_*(\Omega)$ and a unique active index $j_*$,
\begin{equation}
\frac{\dot b_*}{b_*}
=2\operatorname{sgn}(\eta_{j_*})\dot\eta_{j_*}.
\label{eq:minimal_perron_velocity}
\end{equation}

\begin{remark}
\label{rem:tied_active_indices}
When several coordinates attain $\|\eta\|_\infty$, $b_*$ is not differentiable in general.  Directional derivatives or a differential inclusion must replace \eqref{eq:minimal_perron_velocity}.
\end{remark}

On a fixed upper face with tangential first-order motion, the second-order viability condition is
\begin{equation}
\nabla\eta_j\cdot\ddot\Omega
+
\dot\Omega^TD^2\eta_j\dot\Omega
\leq0.
\label{eq:upper_face_acceleration}
\end{equation}
The inequality reverses on a lower face.

\section{Exact modal transport and internally generated amplitudes}
\label{sec:modal_transport}

Assume that every rigid polynomial has distinct roots.  For a root $\lambda$, define the right eigenvector
\begin{equation}
v(\lambda)
=(\lambda^{m-1},\lambda^{m-2},\ldots,\lambda,1)^T.
\label{eq:companion_eigenvector}
\end{equation}
Let
\begin{equation}
V_j=
\left[
v(\lambda_{j,1})\ \cdots\ v(\lambda_{j,m})
\right],
\qquad
\Lambda_j=\operatorname{diag}(\lambda_{j,1},\ldots,\lambda_{j,m}).
\label{eq:modal_matrices}
\end{equation}
If $L_j$ acts for $h_j$ steps and $y_j=V_jc_j$, then
\begin{align}
y_{j+1}
&=V_j\Lambda_j^{h_j}c_j,
\label{eq:window_state_handoff}\\
c_{j+1}
&=T_j\Lambda_j^{h_j}c_j,
\qquad
T_j=V_{j+1}^{-1}V_j.
\label{eq:modal_transport}
\end{align}

\begin{proposition}
\label{prop:lagrange_transport}
The entries of the basis-change matrix are
\begin{equation}
(T_j)_{pq}
=
\prod_{\substack{s=1\\s\neq p}}^{m}
\frac{\lambda_{j,q}-\lambda_{j+1,s}}
{\lambda_{j+1,p}-\lambda_{j+1,s}}.
\label{eq:lagrange_transport_entries}
\end{equation}
\end{proposition}

\begin{proof}
The $p$-th coordinate in the new Vandermonde basis is the cardinal Lagrange polynomial for the nodes $\lambda_{j+1,s}$, evaluated at the old node $\lambda_{j,q}$.  Inverse Vandermonde matrices and Lagrange interpolation are equivalent descriptions of this change of basis.
\end{proof}

For a small root displacement
\begin{equation}
\lambda_{j+1,q}=\lambda_{j,q}+\delta\lambda_{j,q},
\label{eq:small_root_displacement}
\end{equation}
one obtains, for $p\neq q$,
\begin{equation}
(T_j)_{pq}
=-\delta\lambda_{j,q}
\frac{P_j'(\lambda_{j,q})}
{(\lambda_{j,q}-\lambda_{j,p})P_j'(\lambda_{j,p})}
+O(\|\delta\lambda_j\|^2).
\label{eq:offdiagonal_modal_mixing}
\end{equation}
Thus modal mixing becomes large near collisions, consistently with the vanishing of the discriminant.

\begin{proposition}
\label{prop:modal_volume}
The exact modal volume identity is
\begin{align}
\det T_j
&=\frac{\det V_j}{\det V_{j+1}},
\label{eq:transport_determinant}\\
|\det T_j|^2
&=\frac{|\operatorname{Disc}P_j|}
{|\operatorname{Disc}P_{j+1}|}.
\label{eq:discriminant_volume_ratio}
\end{align}
\end{proposition}

\begin{proof}
The first identity follows from $T_j=V_{j+1}^{-1}V_j$.  The square of the absolute Vandermonde determinant equals the absolute discriminant of a monic polynomial with simple roots.
\end{proof}

Iterating \eqref{eq:modal_transport} gives
\begin{equation}
c_J
=T_{J-1}\Lambda_{J-1}^{h_{J-1}}
\cdots
T_0\Lambda_0^{h_0}c_0,
\label{eq:iterated_modal_transport}
\end{equation}
and hence
\begin{equation}
y_J
=V_JT_{J-1}\Lambda_{J-1}^{h_{J-1}}
\cdots
T_0\Lambda_0^{h_0}c_0.
\label{eq:global_state_predictor}
\end{equation}
For the first-coordinate observable,
\begin{equation}
U_J
=e_1^TV_JT_{J-1}\Lambda_{J-1}^{h_{J-1}}
\cdots
T_0\Lambda_0^{h_0}c_0.
\label{eq:global_scalar_predictor}
\end{equation}

\begin{warning}
\label{warn:homogeneous_transport_assumption}
Equations \eqref{eq:iterated_modal_transport}--\eqref{eq:global_scalar_predictor} eliminate independent amplitude fitting only under the homogeneous state-handoff model \eqref{eq:window_state_handoff}.  If each data window is reinitialized, forced, normalized independently, or connected through an additional observation operator, new parameters re-enter and must be counted explicitly.
\end{warning}

\section{Canonical boundary realization and the remaining law}
\label{sec:canonical_boundary}

Whenever $b_*(\Omega)<\infty$, define the canonical minimum-Perron realization
\begin{equation}
L_*(\Omega)=L\bigl(a(b_*(\Omega),\Omega)\bigr).
\label{eq:canonical_boundary_leslie}
\end{equation}
This realization lies on the boundary of the nonnegative cone.  At regular points with unique active strain coordinate, the maps
\begin{equation}
\Omega
\longleftrightarrow
\eta
\longmapsto
b_*(\Omega)
\longmapsto
L_*(\Omega)
\label{eq:canonical_coordinate_chain}
\end{equation}
are locally determined, with the first arrow locally invertible.

\begin{remark}
\label{rem:canonical_not_smooth_everywhere}
The construction \eqref{eq:canonical_boundary_leslie} is canonical but only piecewise smooth because of the maximum norm in $b_*$.  It is also numerically sensitive near the spectral boundary \eqref{eq:critical_region_boundary}.
\end{remark}

The static coordinate chain \eqref{eq:canonical_coordinate_chain} does not by itself determine how a future window is reached.  The next sections derive an exact overlap law directly from the scalar sequence.  For fixed $b$, the special reciprocal structure turns the admissible coefficient differences into a linear problem.  At the first generic stride admitting motion, the entire update is reduced to one scalar innovation.



\section{The fixed-Perron reciprocal affine hull}
\label{sec:reciprocal_affine_hull}

The critical-circle condition has a useful linear precursor.  Fix $b>1$, put $m=2K+1$, and for a coefficient vector $a=(a_1,\ldots,a_m)$ define the synthetic-division coefficients
\begin{align}
q_0(a;b)&=1,
\label{eq:q_zero}\\
q_j(a;b)&=b^j-\sum_{\ell=1}^{j}a_\ell b^{j-\ell},
\qquad 1\leq j\leq 2K.
\label{eq:q_definition}
\end{align}
Then
\begin{equation}
P(z;a)
=(z-b)\left(
 z^{2K}+q_1z^{2K-1}+\cdots+q_{2K}
\right)+P(b;a).
\label{eq:synthetic_division_identity}
\end{equation}

\begin{definition}
\label{def:reciprocal_perron_hull}
The fixed-$b$ reciprocal--Perron affine hull is the set
\begin{equation}
\mathcal C^{\mathrm{rec}}_{b,K}
=
\left\{
a\in\mathbb R^m:
P(b;a)=0,
\quad
q_{2K-j}(a;b)=b^{K-j}q_j(a;b),\ 0\leq j\leq K-1
\right\}.
\label{eq:reciprocal_perron_hull}
\end{equation}
\end{definition}

\begin{proposition}
\label{prop:reciprocal_affine_dimension}
The set \eqref{eq:reciprocal_perron_hull} is an affine subspace of dimension $K$.  Equivalently, there is a rank-$(K+1)$ matrix $R_b$ and a vector $c_b$ such that
\begin{equation}
\mathcal C^{\mathrm{rec}}_{b,K}
=\{a:R_ba=c_b\}.
\label{eq:Rb_affine_form}
\end{equation}
One explicit choice of the $K+1$ affine rows is
\begin{align}
\sum_{\ell=1}^{m}b^{m-\ell}a_\ell
&=b^m,
\label{eq:Rb_perron_row}\\
\sum_{\ell=1}^{2K-j}b^{2K-j-\ell}a_\ell
-
\sum_{\ell=1}^{j}b^{K-\ell}a_\ell
&=b^{2K-j}-b^K,
\qquad 0\leq j\leq K-1.
\label{eq:Rb_reciprocal_rows}
\end{align}
Every rigid critical Leslie coefficient vector with Perron value $b$ belongs to this affine subspace.
\end{proposition}

\begin{proof}
All equations in \eqref{eq:reciprocal_perron_hull} are affine linear in $a$.  The $K$ quantities $q_1,\ldots,q_K$ may be chosen freely.  The reciprocal identities determine $q_{K+1},\ldots,q_{2K}$, after which
\begin{equation}
a_j=bq_{j-1}-q_j,
\qquad 1\leq j\leq2K,
\label{eq:a_from_q}
\end{equation}
and $P(b;a)=0$ determines
\begin{equation}
a_m=bq_{2K}=b^{K+1}.
\label{eq:a_m_from_q}
\end{equation}
Hence the affine dimension is $K$.  For a rigid critical spectrum, $q_j=r^jd_j$ with $r=\sqrt b$, and the palindromy $d_{2K-j}=d_j$ gives exactly the reciprocal equations in \eqref{eq:reciprocal_perron_hull}.
\end{proof}

\begin{remark}
\label{rem:reciprocal_not_full_critical}
Membership in $\mathcal C^{\mathrm{rec}}_{b,K}$ is necessary but not sufficient for all non-Perron roots to lie on $|z|=\sqrt b$.  Put
\begin{equation}
d_j=\frac{q_j}{b^{j/2}}.
\label{eq:d_from_q}
\end{equation}
The reciprocal identities give $d_{2K-j}=d_j$.  Full criticality additionally requires the associated polynomial $H$ in \eqref{eq:H_eta} to have all $K$ roots in $[-2,2]$, together with the Leslie nonnegativity inequalities.
\end{remark}

\section{Overlapping windows and exact mobility rigidity}
\label{sec:overlap_rigidity}

Let $(u_t)$ be any scalar sequence.  In the intended arithmetic application one may take $u_t=M(B_t)$ or a predeclared normalization of those values.  Define the recurrence row
\begin{equation}
h_t=(u_{t-1},u_{t-2},\ldots,u_{t-m})\in\mathbb R^m.
\label{eq:recurrence_row}
\end{equation}
The order-$m$ recurrence window beginning at $n$ consists of the $m$ equations
\begin{equation}
h_ta^{(n)}=u_t,
\qquad
n+m\leq t\leq n+2m-1.
\label{eq:window_recurrence_equations}
\end{equation}
When the corresponding square data matrix is nonsingular, these equations determine $a^{(n)}$ uniquely.

Let a second window begin at $n+s$, with $1\leq s\leq m$.  The two recurrence windows share exactly $m-s$ equations.  Let $O_{n,s}$ be the matrix whose rows are the common $h_t$:
\begin{equation}
O_{n,s}
=
\begin{pmatrix}
h_{n+s+m}\\
\vdots\\
h_{n+2m-1}
\end{pmatrix}
\in\mathbb R^{(m-s)\times m},
\label{eq:overlap_matrix}
\end{equation}
with the empty-matrix convention when $s=m$.

\begin{theorem}[Exact overlap constraint]
\label{thm:exact_overlap_constraint}
Assume that $a^{(n)}$ and $a^{(n+s)}$ both fit their respective windows and both belong to $\mathcal C^{\mathrm{rec}}_{b,K}$ for the same fixed $b$.  Then
\begin{equation}
\delta a_{n,s}:=a^{(n+s)}-a^{(n)}
\in
\ker O_{n,s}\cap\ker R_b.
\label{eq:overlap_kernel_intersection}
\end{equation}
Equivalently, with
\begin{equation}
Z_{n,s}
=
\begin{pmatrix}
O_{n,s}\\
R_b
\end{pmatrix},
\label{eq:stacked_overlap_matrix}
\end{equation}
one has
\begin{equation}
Z_{n,s}\delta a_{n,s}=0.
\label{eq:stacked_overlap_null}
\end{equation}
\end{theorem}

\begin{proof}
Every common recurrence equation is satisfied by both coefficient vectors, so subtraction gives $O_{n,s}\delta a_{n,s}=0$.  Since both vectors satisfy the same affine equations $R_ba=c_b$, subtraction also gives $R_b\delta a_{n,s}=0$.
\end{proof}

\begin{corollary}[Generic mobility threshold]
\label{cor:generic_mobility_threshold}
If the rows of $O_{n,s}$ and $R_b$ are in generic relative position, then
\begin{equation}
\dim\ker Z_{n,s}
=
\max\{0,s-K-1\}.
\label{eq:generic_mobility_dimension}
\end{equation}
Hence:
\begin{align}
s\leq K+1
&\quad\Longrightarrow\quad
\text{the old coefficient vector is locally isolated},
\label{eq:rigid_stride_small}\\
s=K+2
&\quad\Longrightarrow\quad
\text{one generic coefficient direction survives},
\label{eq:critical_stride_one_direction}\\
s>K+2
&\quad\Longrightarrow\quad
s-K-1\text{ generic directions survive}.
\label{eq:large_stride_directions}
\end{align}
\end{corollary}

\begin{proof}
The common-equation matrix contributes $m-s$ generic independent rows and $R_b$ contributes $K+1$ independent rows.  Since $m=2K+1$, the generic nullity is
\begin{equation}
m-\min\{m,m-s+K+1\}
=\max\{0,s-K-1\}.
\end{equation}
\end{proof}

\begin{definition}
\label{def:critical_stride}
For fixed $b$, define the critical stride
\begin{equation}
s_*=K+2.
\label{eq:critical_stride}
\end{equation}
It is the first stride at which a nonzero fixed-$b$ reciprocal--Perron motion exists generically.
\end{definition}

\begin{remark}
\label{rem:small_stride_degeneracy}
For $s\leq K+1$, nontrivial motion is possible only when $Z_{n,s}$ loses rank.  Thus exceptional motion is detected by the vanishing of maximal minors of $Z_{n,s}$.  In particular, any single Hankel-type minor obtained from the fixed identity $a_m=b^{K+1}$ is only one necessary degeneracy condition, not a sufficient condition for full critical motion.
\end{remark}

\begin{remark}
\label{rem:moving_b_threshold}
If $b$ is allowed to vary, the regular rigid critical manifold has dimension $K+1$.  A transversality count gives local intersection dimension
\begin{equation}
\max\{0,s-K\},
\label{eq:moving_b_local_dimension}
\end{equation}
so the first generic one-dimensional mobility occurs at $s=K+1$.  Unlike the fixed-$b$ result, this statement is local and nonlinear because coefficient differences for distinct Perron values do not lie in one fixed affine subspace.
\end{remark}

\section{Scalar innovation at the critical stride}
\label{sec:scalar_innovation}

Set $s=s_*=K+2$ and assume
\begin{equation}
\operatorname{rank}Z_{n,s_*}=m-1.
\label{eq:critical_stride_rank}
\end{equation}
Choose a nonzero vector $v_n$ spanning the one-dimensional nullspace,
\begin{equation}
\ker Z_{n,s_*}=\operatorname{span}\{v_n\}.
\label{eq:critical_stride_direction}
\end{equation}

\begin{theorem}[Scalarization of the window update]
\label{thm:scalarization}
Under \eqref{eq:critical_stride_rank}, every next fixed-$b$ reciprocal--Perron coefficient vector that agrees with the old window on all common equations has the form
\begin{equation}
a^{(n+s_*)}=a^{(n)}+\tau_n v_n
\label{eq:scalar_coefficient_update}
\end{equation}
for a unique scalar $\tau_n$.
\end{theorem}

\begin{proof}
The difference belongs to the one-dimensional nullspace in \eqref{eq:critical_stride_direction}.
\end{proof}

The next window contributes $s_*=K+2$ genuinely new recurrence equations.  The first occurs at
\begin{equation}
t_0=n+2m.
\label{eq:first_new_time}
\end{equation}
Its recurrence row $h_{t_0}$ uses only values already known before $u_{t_0}$ is observed.

\begin{theorem}[One innovation sample determines the next Leslie]
\label{thm:one_sample_determines_update}
Assume
\begin{equation}
h_{t_0}v_n\neq0.
\label{eq:innovation_observability}
\end{equation}
Then the first new scalar value determines the update parameter exactly:
\begin{equation}
\boxed{
\tau_n
=
\frac{u_{t_0}-h_{t_0}a^{(n)}}{h_{t_0}v_n}.
}
\label{eq:tau_from_first_new_value}
\end{equation}
Consequently the entire next coefficient vector is fixed by \eqref{eq:scalar_coefficient_update}.
\end{theorem}

\begin{proof}
Substitute \eqref{eq:scalar_coefficient_update} into the recurrence equation at $t_0$ and solve the resulting scalar equation for $\tau_n$.
\end{proof}

\begin{corollary}[Causal $K+1$-value forecast]
\label{cor:causal_Kplus1_forecast}
After observing only $u_{t_0}$, define $a^+=a^{(n)}+\tau_nv_n$ using \eqref{eq:tau_from_first_new_value}.  Then recursively define
\begin{equation}
\widehat u_t
=
\sum_{\ell=1}^{m}a^+_\ell\widehat u_{t-\ell},
\qquad
t_0+1\leq t\leq t_0+K+1,
\label{eq:causal_recursive_forecast}
\end{equation}
using the observed historical values for indices below $t_0$ and earlier forecasts above $t_0$.  If the next window is exactly generated by a fixed-$b$ rigid critical Leslie coefficient vector, then
\begin{equation}
\widehat u_t=u_t,
\qquad
t_0+1\leq t\leq t_0+K+1.
\label{eq:exact_Kplus1_forecast}
\end{equation}
\end{corollary}

\begin{proof}
The first new observation determines the exact next coefficient vector.  The remaining $K+1$ new equations of that window use the same coefficient vector.  Induction on $t$ in \eqref{eq:causal_recursive_forecast} gives the result.
\end{proof}

\begin{proposition}[Rank-one block diagnostic]
\label{prop:rank_one_block_diagnostic}
Let $G_n$ contain the $K+2$ new recurrence rows of the next window, evaluated from the observed sequence, and let $y_n^+$ be their right-hand side.  Define
\begin{equation}
e_n=y_n^+-G_na^{(n)},
\qquad
g_n=G_nv_n.
\label{eq:block_residual_vectors}
\end{equation}
Then exact critical-stride transport implies
\begin{equation}
e_n=\tau_ng_n.
\label{eq:block_collinearity}
\end{equation}
Equivalently,
\begin{equation}
(e_n)_i(g_n)_j-(e_n)_j(g_n)_i=0
\qquad
\text{for all }i,j.
\label{eq:block_minor_identities}
\end{equation}
Thus a block of $K+2$ new equations carries only one scalar innovation after the old state and overlap direction are fixed.
\end{proposition}

\section{Admissible innovation interval and boundary locking}
\label{sec:innovation_interval}

The scalar parameter in \eqref{eq:scalar_coefficient_update} cannot take arbitrary values if nonnegativity and the critical-circle root condition are retained.

\begin{definition}
\label{def:positive_tau_interval}
For a current nonnegative coefficient vector $a$ and direction $v$, define the interval of coefficient positivity
\begin{equation}
I_{\mathrm{pos}}(a,v)
=
\left[
\max_{v_\ell>0}\left(-\frac{a_\ell}{v_\ell}\right),
\ 
\min_{v_\ell<0}\left(-\frac{a_\ell}{v_\ell}\right)
\right],
\label{eq:positive_tau_interval}
\end{equation}
with the usual $-\infty$ and $+\infty$ conventions when one of the indexing sets is empty.
\end{definition}

For $a(\tau)=a+\tau v$, define $q_j(\tau)$ by \eqref{eq:q_definition}, $d_j(\tau)=q_j(\tau)/b^{j/2}$, and the corresponding reciprocal reduction $H_\tau$.  Let $I_{\mathrm{spec}}$ be the connected component containing $0$ of the set of $\tau$ for which $H_\tau$ has $K$ simple roots in $(-2,2)$.  At regular points, an endpoint of this component can occur only when
\begin{equation}
\operatorname{Disc}H_\tau=0
\qquad\text{or}\qquad
H_\tau(-2)=0.
\label{eq:tau_spectral_boundaries}
\end{equation}
Define the full local innovation interval
\begin{equation}
I_{\mathrm{adm}}=I_{\mathrm{pos}}\cap I_{\mathrm{spec}}.
\label{eq:admissible_tau_interval}
\end{equation}

\begin{corollary}[Pre-observation admissible range]
\label{cor:first_value_admissible_range}
Before $u_{t_0}$ is observed, the moving rigid Leslie hypothesis requires
\begin{equation}
\boxed{
 u_{t_0}
 \in
 h_{t_0}a^{(n)}
 +(h_{t_0}v_n)I_{\mathrm{adm}}.
}
\label{eq:first_value_allowed_interval}
\end{equation}
A value outside this interval falsifies the local fixed-$b$ critical transition before any later value of the new block is used.
\end{corollary}

\begin{proposition}[Boundary locking criterion]
\label{prop:boundary_locking}
Let
\begin{equation}
\mathcal Z(a)=\{\ell:a_\ell=0\}
\label{eq:active_zero_set}
\end{equation}
be the active coefficient faces.  A nonzero one-sided coefficient motion $a+\tau v$ preserving nonnegativity exists locally only if the numbers $v_\ell$, $\ell\in\mathcal Z(a)$, do not contain both positive and negative values.  If both signs occur, then $\tau=0$ is the only sufficiently small feasible update.
\end{proposition}

\begin{proof}
For each active coordinate the inequality $a_\ell+\tau v_\ell\geq0$ reduces to $\tau v_\ell\geq0$.  A common nonzero sign of $\tau$ exists exactly when the active $v_\ell$ do not impose conflicting signs.
\end{proof}

\begin{remark}
\label{rem:minimum_perron_locking}
The canonical minimum-Perron realization $L_*(\Omega)$ always has at least one active coefficient face.  It is therefore dynamically more restrictive than an interior choice $b>b_*(\Omega)$.  A canonical boundary trajectory may lock whenever the critical-stride direction intersects active faces with incompatible signs.
\end{remark}

\section{Exact scalar-driven strain transport}
\label{sec:scalar_strain_transport}

For the update $a^+=a+\tau v$, define
\begin{equation}
\psi_j(v;b)
=
\sum_{\ell=1}^{j}v_\ell b^{j-\ell},
\qquad
\psi_0=0.
\label{eq:psi_definition}
\end{equation}
Equation \eqref{eq:q_definition} gives
\begin{equation}
q_j(a+\tau v;b)=q_j(a;b)-\tau\psi_j(v;b).
\label{eq:q_scalar_update}
\end{equation}

\begin{theorem}[Exact strain update]
\label{thm:exact_strain_update}
Whenever all logarithms are defined,
\begin{equation}
\boxed{
\eta_j^+-\eta_j
=
\log\left(1-\tau\frac{\psi_j}{q_j}\right)
-
\log\left(1-\tau\frac{\psi_{j-1}}{q_{j-1}}\right),
\qquad 1\leq j\leq K.
}
\label{eq:exact_eta_scalar_update}
\end{equation}
Hence, at the critical stride, all $K$ strain coordinates are driven by the same scalar innovation $\tau$.
\end{theorem}

\begin{proof}
Use $d_j=q_j/b^{j/2}$ in the definition $\eta_j=\log(d_j/d_{j-1})$ and substitute \eqref{eq:q_scalar_update}.
\end{proof}

Differentiating with respect to the scalar innovation gives the exact line-flow equation
\begin{equation}
\frac{d\eta_j}{d\tau}
=
-\frac{\psi_j}{q_j-\tau\psi_j}
+
\frac{\psi_{j-1}}{q_{j-1}-\tau\psi_{j-1}}.
\label{eq:eta_tau_flow}
\end{equation}
For small $\tau$,
\begin{equation}
\eta_j^+-\eta_j
=
-\tau
\left(
\frac{\psi_j}{q_j}
-
\frac{\psi_{j-1}}{q_{j-1}}
\right)
+O(\tau^2).
\label{eq:eta_small_tau}
\end{equation}
Thus the $K$-dimensional strain increment is generically confined to one first-order direction at the critical stride.

\section{Frequency-free determinant identities}
\label{sec:frequency_free_determinants}

The reciprocal--Perron test can be imposed directly on a data window without first estimating any phases.  Let
\begin{equation}
\mathsf A_n
=
\begin{pmatrix}
h_{n+m}\\
\vdots\\
h_{n+2m-1}
\end{pmatrix},
\qquad
\mathsf g_n
=
\begin{pmatrix}
u_{n+m}\\
\vdots\\
u_{n+2m-1}
\end{pmatrix}.
\label{eq:data_recurrence_matrix}
\end{equation}
Write the affine reciprocal constraints in \eqref{eq:Rb_affine_form} row by row as
\begin{equation}
\rho_\alpha(b)^Ta=\kappa_\alpha(b),
\qquad
0\leq\alpha\leq K.
\label{eq:affine_constraint_rows}
\end{equation}
For each row define
\begin{equation}
\Delta_{\alpha,n}(b)
=
\det
\begin{pmatrix}
\mathsf A_n & \mathsf g_n\\
\rho_\alpha(b)^T & \kappa_\alpha(b)
\end{pmatrix},
\label{eq:frequency_free_determinant}
\end{equation}

\begin{proposition}
\label{prop:frequency_free_test}
If $\det\mathsf A_n\neq0$, then the unique order-$m$ recurrence coefficient vector of the window belongs to $\mathcal C^{\mathrm{rec}}_{b,K}$ if and only if
\begin{equation}
\Delta_{\alpha,n}(b)=0,
\qquad 0\leq\alpha\leq K.
\label{eq:all_frequency_free_determinants}
\end{equation}
Full rigid criticality is obtained by adding the nonnegativity and root-location tests of Sections~\ref{sec:eta_coordinates}--\ref{sec:memory_critical_region}.
\end{proposition}

\begin{proof}
Since $a=\mathsf A_n^{-1}\mathsf g_n$, the Schur-complement identity gives
\begin{equation}
\Delta_{\alpha,n}(b)
=
\det(\mathsf A_n)
\left(
\kappa_\alpha(b)-\rho_\alpha(b)^T\mathsf A_n^{-1}\mathsf g_n
\right).
\label{eq:schur_determinant_identity}
\end{equation}
Thus the determinants vanish exactly when all affine reciprocal--Perron constraints hold.
\end{proof}

\begin{remark}
\label{rem:unknown_b_elimination}
If $b$ is not fixed in advance, the functions $\Delta_{\alpha,n}(b)$ are explicit polynomials in $b$ whose coefficients depend only on the data window.  A necessary condition for a rigid critical realization with some $b>1$ is that these $K+1$ polynomials possess a common real root $b>1$.  This can be tested without frequencies by polynomial gcd or resultant/subresultant elimination, followed by the root-location and positivity conditions.
\end{remark}

\section{Closed form for one mode pair}
\label{sec:K1_closed_form}

For $K=1$ one has $m=3$, and the entire construction becomes explicit.  Write $\alpha=a_1$.  The fixed-$b$ reciprocal--Perron constraints give
\begin{equation}
\boxed{
 a_1=\alpha,
 \qquad
 a_2=b^2-b-b\alpha,
 \qquad
 a_3=b^2.
}
\label{eq:K1_coefficients}
\end{equation}
The quotient polynomial is
\begin{equation}
Q(z)=z^2+(b-\alpha)z+b.
\label{eq:K1_quotient}
\end{equation}
Its two roots lie on $|z|=\sqrt b$, and the Leslie coefficients are nonnegative, if and only if
\begin{equation}
\boxed{
\max\{0,b-2\sqrt b\}
\leq\alpha\leq b-1.
}
\label{eq:K1_alpha_interval}
\end{equation}
Endpoint equalities include coefficient or spectral degeneracies; strict interior conditions give the corresponding nondegenerate case.

Define
\begin{align}
X_t(b)&=u_{t-1}-bu_{t-2},
\label{eq:K1_X}\\
Y_t(b)&=u_t-(b^2-b)u_{t-2}-b^2u_{t-3}.
\label{eq:K1_Y}
\end{align}
Then the critical recurrence is simply
\begin{equation}
Y_t(b)=\alpha X_t(b).
\label{eq:K1_rank_one_relation}
\end{equation}
Thus, if $(X_t,X_{t+1},X_{t+2})$ is not identically zero, three consecutive recurrence equations admit a common fixed-$b$ reciprocal coefficient vector if and only if all pairwise minors vanish:
\begin{equation}
X_iY_j-X_jY_i=0,
\qquad
t\leq i<j\leq t+2.
\label{eq:K1_data_identities}
\end{equation}
If $X_t=X_{t+1}=X_{t+2}=0$, existence instead requires $Y_t=Y_{t+1}=Y_{t+2}=0$.  Full rigid criticality additionally requires the inferred $\alpha$ to satisfy \eqref{eq:K1_alpha_interval}.

For $K=1$, the critical stride is $s_*=3$ and the fixed-$b$ affine direction is
\begin{equation}
v=(1,-b,0)^T.
\label{eq:K1_direction}
\end{equation}
Hence
\begin{equation}
\alpha^+=\alpha+\tau,
\qquad
a_2^+=a_2-b\tau,
\qquad
a_3^+=b^2.
\label{eq:K1_scalar_update}
\end{equation}
The first new recurrence equation gives
\begin{equation}
\tau
=
\frac{u_{t_0}-h_{t_0}a}{X_{t_0}(b)},
\label{eq:K1_tau}
\end{equation}
provided $X_{t_0}(b)\neq0$, and the next two values are then fixed recursively.  The exact admissible scalar interval is
\begin{equation}
\boxed{
\max\{0,b-2\sqrt b\}-\alpha
\leq\tau
\leq
b-1-\alpha.
}
\label{eq:K1_tau_interval}
\end{equation}
This gives a completely frequency-free three-step falsification and two-value forecasting rule for the one-pair model.

\section{Irreducible scalar forcing and the endpoint of the Leslie geometry}
\label{sec:irreducible_scalar_forcing}

\begin{proposition}
\label{prop:irreducible_scalar_forcing}
At a regular interior fixed-$b$ point and the critical stride $s_*=K+2$, the Leslie geometry cannot determine $\tau_n$ from the old window and the common equations alone.  In a neighborhood of $\tau=0$, distinct values of $\tau$ produce distinct reciprocal--Perron coefficient vectors satisfying all those old and common constraints.  Full criticality persists for a nontrivial interval whenever the current point lies in the interior of the fixed-$b$ critical region.
\end{proposition}

\begin{proof}
The rank assumption \eqref{eq:critical_stride_rank} leaves exactly the affine line \eqref{eq:scalar_coefficient_update}.  Within the reciprocal affine hull, coefficient positivity and the condition that the reduced polynomial have simple roots in $(-2,2)$ are open.  Hence sufficiently small values of $\tau$ remain admissible.
\end{proof}

\begin{remark}
\label{rem:true_remaining_problem}
This identifies the exact limit of the moving-Leslie mechanism.  The geometry compresses a block of $K+2$ genuinely new recurrence equations to one scalar innovation and then transports all $K$ strain coordinates through \eqref{eq:exact_eta_scalar_update}.  It does not supply a law for that scalar innovation.  A zero-latency forecast of the first new value therefore requires an additional arithmetic statement about the sequence $\tau_n$.  Any such law is external to Leslie geometry and is the remaining place where genuinely Möbius-specific information would have to enter.
\end{remark}

\section{Exact first-order spectral transport}
\label{sec:first_order_transport}

Let $a(t)$ be a differentiable coefficient path and define
\begin{equation}
P(z,t)=z^m-\sum_{\ell=1}^{m}a_\ell(t)z^{m-\ell}.
\label{eq:time_polynomial}
\end{equation}
Let $\lambda_k(t)$ be a simple root, so that
\begin{equation}
P(\lambda_k(t),t)=0,
\qquad
P_z(\lambda_k(t),t)\neq0.
\label{eq:simple_root_hypothesis}
\end{equation}
Differentiability of simple eigenvalue branches is standard in matrix perturbation theory \cite{Yener2023}.

\begin{proposition}
\label{prop:first_derivative}
Under \eqref{eq:simple_root_hypothesis},
\begin{equation}
\dot\lambda_k
=
\frac{
\displaystyle
\sum_{\ell=1}^{m}
\dot a_\ell\lambda_k^{m-\ell}
}{P_z(\lambda_k,t)}.
\label{eq:eigenvalue_velocity}
\end{equation}
\end{proposition}

\begin{proof}
Differentiating $P(\lambda_k(t),t)=0$ gives
\begin{equation}
P_z(\lambda_k,t)\dot\lambda_k
+P_t(\lambda_k,t)=0.
\label{eq:implicit_derivative}
\end{equation}
Since
\begin{equation}
P_t(z,t)
=-\sum_{\ell=1}^{m}\dot a_\ell z^{m-\ell},
\label{eq:polynomial_time_derivative}
\end{equation}
substitution into \eqref{eq:implicit_derivative} yields \eqref{eq:eigenvalue_velocity}.
\end{proof}

Suppose now that
\begin{equation}
\lambda_k(t)=\sqrt{b(t)}e^{i\omega_k(t)}.
\label{eq:critical_root_parameterization}
\end{equation}
Then
\begin{equation}
\frac{\dot\lambda_k}{\lambda_k}
=
\frac{\dot b}{2b}+i\dot\omega_k.
\label{eq:logarithmic_root_velocity}
\end{equation}
Define
\begin{equation}
J_{k\ell}(L)
=
\frac{\lambda_k^{m-\ell-1}}{P_z(\lambda_k)}.
\label{eq:spectral_transport_matrix}
\end{equation}

\begin{corollary}
\label{cor:phase_and_radius_velocity}
For every simple prescribed branch,
\begin{align}
\dot\omega_k
&=
\Im\left(
\sum_{\ell=1}^{m}J_{k\ell}(L)\dot a_\ell
\right),
\label{eq:angular_velocity}\\
\frac{\dot b}{2b}
&=
\Re\left(
\sum_{\ell=1}^{m}J_{k\ell}(L)\dot a_\ell
\right).
\label{eq:radial_velocity}
\end{align}
\end{corollary}

\begin{warning}
\label{warn:radial_compatibility}
Equation \eqref{eq:radial_velocity} must give the same scalar $\dot b/(2b)$ for every $k$.  Thus the imaginary-part formula alone is not a complete admissibility test.  The real-part compatibility equations are additional local constraints.
\end{warning}

\section{The rigid critical manifold}
\label{sec:critical_manifold}

Let
\begin{equation}
\mathcal M_K
=
\left\{
a(b,\Omega):
 b>1,\ \Omega\in\mathbb T^K,
\ a(b,\Omega)\geq0
\right\}
\subseteq\mathbb R^{2K+1}.
\label{eq:rigid_critical_set}
\end{equation}

\begin{proposition}
\label{prop:generic_dimension}
At a point where the coefficient map $(b,\Omega)\mapsto a(b,\Omega)$ has rank $K+1$ and the prescribed roots are distinct, $\mathcal M_K$ is locally a smooth manifold of dimension $K+1$ and codimension $K$ in $\mathbb R^{2K+1}$.
\end{proposition}

\begin{proof}
The parameter space has dimension $K+1$.  Under the stated full-rank hypothesis, the constant-rank theorem gives a local image manifold of dimension $K+1$.  The ambient dimension is $2K+1$, so the codimension is $K$.
\end{proof}

Consequently, locally there are $K$ independent constraints
\begin{equation}
F_r(a)=0,
\qquad 1\leq r\leq K,
\label{eq:implicit_constraints}
\end{equation}
and every differentiable path in the regular part satisfies
\begin{equation}
\nabla F_r(a)\cdot\dot a=0,
\qquad 1\leq r\leq K.
\label{eq:tangent_equalities}
\end{equation}

\begin{warning}
\label{warn:singular_manifold}
The dimension statement is generic, not universal.  Repeated roots, collisions of phases, loss of rank, and intersections with several active coefficient faces can create singular or lower-dimensional strata.
\end{warning}

\section{Positivity and the corrected spectral velocity cone}
\label{sec:velocity_cone}

For $a\in\mathbb R_+^m$, the tangent cone of the nonnegative orthant is
\begin{equation}
T_a\mathbb R_+^m
=
\left\{
v\in\mathbb R^m:
 v_\ell\geq0\text{ whenever }a_\ell=0
\right\}.
\label{eq:orthant_tangent_cone}
\end{equation}
Tangent cones generalize tangent spaces at nonsmooth boundary points of convex sets \cite{Wright2019}.

The velocity must satisfy three types of conditions simultaneously:
\begin{enumerate}
\item orthant feasibility, $v\in T_a\mathbb R_+^m$;
\item tangency to the rigid critical manifold, $DF(a)v=0$;
\item common radial motion across all prescribed root pairs.
\end{enumerate}

Define the common-radial constraint operator by
\begin{equation}
\mathcal R_Lv
=
\left(
\Re(J_{2\cdot}(L)v)-\Re(J_{1\cdot}(L)v),
\ldots,
\Re(J_{K\cdot}(L)v)-\Re(J_{1\cdot}(L)v)
\right)^T.
\label{eq:common_radial_operator}
\end{equation}
At a regular point, define the admissible coefficient-velocity cone
\begin{equation}
\mathcal K(L)
=
T_a\mathbb R_+^m
\cap\ker DF(a)
\cap\ker\mathcal R_L.
\label{eq:admissible_coefficient_cone}
\end{equation}

\begin{definition}
\label{def:spectral_velocity_cone}
The rigid positive spectral velocity cone is
\begin{equation}
\mathcal V(L)
=
\left\{
\Im J(L)v:
 v\in\mathcal K(L)
\right\}.
\label{eq:spectral_velocity_cone}
\end{equation}
\end{definition}

\begin{proposition}
\label{prop:necessary_velocity_condition}
Every differentiable rigid positive Leslie path with simple prescribed roots satisfies
\begin{equation}
\dot\Omega(t)\in\mathcal V(L(t)).
\label{eq:velocity_necessary_condition}
\end{equation}
\end{proposition}

\begin{proof}
The coefficient velocity lies in the orthant tangent cone, is tangent to the rigid critical manifold, and satisfies the common-radial identities.  Equation \eqref{eq:angular_velocity} then places its phase velocity in the image \eqref{eq:spectral_velocity_cone}.
\end{proof}

\begin{remark}
\label{rem:necessary_not_sufficient}
Condition \eqref{eq:velocity_necessary_condition} is first-order and necessary.  At singular points or for finite window increments it is not sufficient for the existence of a global path.
\end{remark}

\section{Second-order transport}
\label{sec:second_order_transport}

\begin{proposition}
\label{prop:second_derivative}
If $a(t)$ is twice differentiable and \eqref{eq:simple_root_hypothesis} holds, then
\begin{equation}
\ddot\lambda_k
=-
\frac{
P_{zz}\dot\lambda_k^2
+2P_{zt}\dot\lambda_k
+P_{tt}
}{P_z},
\label{eq:eigenvalue_acceleration}
\end{equation}
where all polynomial derivatives are evaluated at $(\lambda_k(t),t)$ and
\begin{equation}
P_{tt}(z,t)
=-\sum_{\ell=1}^{m}\ddot a_\ell z^{m-\ell}.
\label{eq:second_time_derivative}
\end{equation}
\end{proposition}

\begin{proof}
Differentiate \eqref{eq:implicit_derivative} once more and collect the chain-rule terms.  Division by $P_z\neq0$ gives \eqref{eq:eigenvalue_acceleration}.
\end{proof}

For the parameterization \eqref{eq:critical_root_parameterization},
\begin{equation}
\frac{\ddot\lambda_k}{\lambda_k}
=
\frac{\ddot b}{2b}
-\frac{\dot b^2}{4b^2}
-\dot\omega_k^2
+i\left(
\ddot\omega_k+\frac{\dot b}{b}\dot\omega_k
\right).
\label{eq:angular_acceleration_identity}
\end{equation}
Thus a smooth Leslie path constrains not only $\Omega$ and $\dot\Omega$, but also the angular acceleration after the radial terms are removed.

\begin{definition}
\label{def:spectral_curvature}
At a regular path point, the residual obtained by comparing the observed $\ddot\Omega$ with the imaginary part implied by \eqref{eq:eigenvalue_acceleration} is called the local Leslie spectral-curvature residual.
\end{definition}

\begin{warning}
\label{warn:curvature_not_intrinsic}
A numerical curvature residual depends on the chosen interpolation of discrete windows and on the time scale assigned to the window index.  It is not intrinsic until those choices are fixed before testing.
\end{warning}

\section{Discrete windows and finite increments}
\label{sec:discrete_windows}

The observations are discrete, not a differentiable curve.  Let
\begin{equation}
\Delta a_j=a_{j+1}-a_j,
\qquad
\Delta\Omega_j=\operatorname{unwrap}(\Omega_{j+1}-\Omega_j).
\label{eq:window_increments}
\end{equation}
The differential model suggests the residual
\begin{equation}
r_j
=
\Delta\Omega_j-
\Im J(L_j)\Delta a_j.
\label{eq:first_order_window_residual}
\end{equation}

\begin{warning}
\label{warn:finite_increment}
Equation \eqref{eq:first_order_window_residual} is only a first-order approximation.  A valid finite-step test should also check that the segment or reconstructed path remains in the admissible set and should report an error scale of order $\|\Delta a_j\|^2$ away from singularities.
\end{warning}

A stronger discrete feasibility test is to solve for a path $a_j(s)$, $0\leq s\leq1$, satisfying
\begin{equation}
a_j(0)=a_j,
\qquad
a_j(1)=a_{j+1},
\qquad
a_j(s)\in\mathcal M_K\cap\mathbb R_+^m.
\label{eq:path_between_windows}
\end{equation}
Existence of \eqref{eq:path_between_windows} is stronger than endpoint feasibility.

\section{A genuine forecasting protocol}
\label{sec:forecasting_protocol}

\begin{experiment}
\label{exp:window_forecast}
Choose before examining the test windows:
\begin{enumerate}
\item a fixed base $b$ or a predeclared rule for $b_j$;
\item a fixed number $K$ of mode pairs and hence $m=2K+1$;
\item consecutive, nonoverlapping or explicitly overlapping logarithmic windows;
\item a frequency estimator and a deterministic mode-matching rule;
\item a coefficient-transport law trained only on past windows;
\item a rule for handling windows outside $\mathcal A_{b,K}$;
\item a horizon and evaluation metric.
\end{enumerate}
For each training endpoint $J$, infer $L_0,\ldots,L_J$, predict $\widehat a_{J+1}$ without using $W_{J+1}$, project to the admissible rigid set if that projection was predeclared, compute $\widehat\Omega_{J+1}$ from the roots, and only then evaluate the prediction on $W_{J+1}$.
\end{experiment}

The information flow must be
\begin{equation}
M|_{W_0,\ldots,W_J}
\longrightarrow
L_0,\ldots,L_J
\longrightarrow
\widehat L_{J+1}
\longrightarrow
\widehat\Omega_{J+1}
\longrightarrow
\widehat M|_{W_{J+1}}.
\label{eq:forecast_information_flow}
\end{equation}

\begin{warning}
\label{warn:law_required}
The identities of Sections~\ref{sec:first_order_transport}--\ref{sec:second_order_transport} do not themselves specify a forecasting law for $\dot a$ or $\ddot a$.  A predictor requires a predeclared dynamical model for coefficient evolution.  Merely estimating $L_{J+1}$ from the future window would return to independent fitting.
\end{warning}


\section{Critical-stride forecasting protocol}
\label{sec:critical_stride_protocol}

\begin{experiment}
\label{exp:critical_stride_forecast}
Fix $K$, $m=2K+1$, a Perron value $b>1$, and the critical stride $s_*=K+2$ before examining the test region.  For each training window:
\begin{enumerate}
\item solve the $m$ recurrence equations \eqref{eq:window_recurrence_equations} directly from the past scalar sequence;
\item reject the window unless the resulting coefficients satisfy the fixed-$b$ reciprocal determinants \eqref{eq:all_frequency_free_determinants}, nonnegativity, and the critical-circle root-location test;
\item construct $Z_{n,s_*}$ using only common equations and $R_b$;
\item require $\operatorname{rank}Z_{n,s_*}=m-1$ and compute its null direction $v_n$;
\item compute the pre-observation interval \eqref{eq:first_value_allowed_interval};
\item reveal only the first new scalar $u_{t_0}$ and compute $\tau_n$ from \eqref{eq:tau_from_first_new_value};
\item without recalibration, predict the remaining $K+1$ new values using \eqref{eq:causal_recursive_forecast};
\item record the $K+1$ forecast errors, the scalar innovation $\tau_n$, the strain update \eqref{eq:exact_eta_scalar_update}, and whether $\tau_n\in I_{\mathrm{adm}}$.
\end{enumerate}
\end{experiment}

The information flow in this strongest fixed-$b$ test is therefore
\begin{equation}
\text{past block}
\longrightarrow
(a^{(n)},v_n,I_{\mathrm{adm}})
\longrightarrow
u_{t_0}
\longrightarrow
\tau_n
\longrightarrow
\widehat u_{t_0+1},\ldots,\widehat u_{t_0+K+1},
\label{eq:critical_stride_information_flow}
\end{equation}
where $u_{t_0}$ is the single newly revealed scalar value; no later test value is used to update the coefficients.

\begin{remark}
\label{rem:innovation_sequence_next_target}
If the $K+1$-value forecasts succeed repeatedly, the remaining arithmetic object is the scalar innovation sequence $(\tau_n)$.  A further predictive law for $\tau_n$ would upgrade the one-sample-delayed forecast to a genuine zero-latency block forecast.  Leslie geometry alone does not provide such a law by Proposition~\ref{prop:irreducible_scalar_forcing}.
\end{remark}

\section{Three nested models}
\label{sec:nested_models}

\begin{definition}
\label{def:models_abc}
The comparison should separate:
\begin{enumerate}
\item free frequencies fitted independently in every window;
\item rigid nonnegative Leslie matrices fitted independently in every window;
\item one dynamically transported rigid nonnegative Leslie trajectory, forecast from past windows only.
\end{enumerate}
\end{definition}

Model B tests static realizability.  Model C tests whether a low-complexity trajectory inside the constrained geometry has forecasting content.  Model C should be compared not only with Models A and B, but also with persistence, autoregressive, phase-randomized, and support-preserving null models.

\section{What the previous forty runs do and do not establish}
\label{sec:forty_runs}

\begin{experiment}
\label{exp:previous_runs}
The reported forty runs varied the sampling base, training fraction, and requested number of modes.  They established numerical out-of-sample prediction by free frequencies and static positive spectral completion with auxiliary eigenvalues for those forty configurations.
\end{experiment}

\begin{proposition}
\label{prop:runs_not_trajectory}
The forty configurations do not constitute a temporal Leslie trajectory and cannot test \eqref{eq:velocity_necessary_condition} or the forecast protocol \eqref{eq:forecast_information_flow}.
\end{proposition}

\begin{proof}
A temporal test requires consecutively ordered windows under fixed or predeclared modelling choices.  The forty runs index alternative parameter configurations rather than successive time windows of one fixed model.
\end{proof}

\section{Verified corrections to the proposed formulation}
\label{sec:verified_corrections}

\begin{proposition}
\label{prop:verified_formulae}
The first derivative \eqref{eq:eigenvalue_velocity}, the angular and radial decompositions \eqref{eq:angular_velocity}--\eqref{eq:radial_velocity}, and the second derivative \eqref{eq:eigenvalue_acceleration} are algebraically correct under the simple-root hypothesis.
\end{proposition}

\begin{remark}
\label{rem:numerical_check}
An independent finite-difference calculation was performed for a degree-five companion polynomial with randomly chosen coefficient velocities.  The maximum discrepancy between the finite-difference root velocity and \eqref{eq:eigenvalue_velocity} was below $9\times10^{-9}$ for a step of $10^{-7}$.  This is a numerical consistency check, not a substitute for Proposition~\ref{prop:first_derivative}.
\end{remark}

\begin{remark}
\label{rem:corrections}
Three corrections are essential:
\begin{enumerate}
\item the boundary relation is the inclusion \eqref{eq:admissible_boundary}, not automatic equality;
\item the spectral velocity cone must include tangency to the rigid manifold and common radial compatibility, as in \eqref{eq:admissible_coefficient_cone};
\item the dimension and codimension counts are generic and require a full-rank hypothesis.
\end{enumerate}
\end{remark}


\section{Numerical obstruction for the rigid $K=7$ core}
\label{sec:K7_rigid_obstruction}

The rigid feasibility criterion of Corollary~\ref{cor:ratio_test} can be evaluated even when the logarithmic coordinates $\eta_j$ do not exist.  For $r=\sqrt b$ define
\begin{equation}
\delta_j(b,\Omega)
=
\min\left\{
 r d_{j-1}-d_j,
 r d_j-d_{j-1}
\right\},
\qquad 1\leq j\leq K,
\label{eq:paired_rigid_margin}
\end{equation}
and
\begin{equation}
\Delta_K(b,\Omega)
=
\min_{1\leq j\leq K}\delta_j(b,\Omega).
\label{eq:global_rigid_margin}
\end{equation}
Because the factors multiplying the two quantities in \eqref{eq:paired_rigid_margin} in the paired Leslie coefficients are positive, one has
\begin{equation}
\Delta_K(b,\Omega)\geq0
\quad\Longleftrightarrow\quad
P_{b,\Omega}
\text{ is the characteristic polynomial of a nonnegative rigid Leslie matrix}.
\label{eq:delta_exact_feasibility}
\end{equation}
Unlike $b_*(\Omega)$, the margin \eqref{eq:global_rigid_margin} remains defined when one or more $d_j$ are nonpositive.

\begin{experiment}
\label{exp:K7_rigid_obstruction}
A grid of $241$ values of $b$ in $[1.01,1.25]$ was examined with $K=7$ and the predeclared frequency estimator.  The estimator returned all seven frequencies at eight grid points.  For every one of these eight evaluable configurations,
\begin{equation}
\Delta_7(b,\Omega)<0.
\label{eq:K7_all_negative}
\end{equation}
The largest observed margin was
\begin{equation}
\boxed{
\Delta_{\max}=-1.228610846762
}
\qquad\text{at}\qquad
\boxed{b=1.178}.
\label{eq:K7_delta_max}
\end{equation}
Thus none of the eight evaluable critical spectra is the complete spectrum of a nonnegative Leslie matrix of rigid dimension $15$.

Seven of the eight configurations had $d_1,\ldots,d_7>0$ and therefore admitted the logarithmic coordinates $\eta$ and the finite quantity $b_*(\Omega)$.  In all seven cases $b_*(\Omega)>b$.  The eighth configuration had all seven frequencies but at least one nonpositive $d_j$, so $b_*$ was not finite although $\Delta_7$ remained evaluable.

The index attaining the most negative paired margin had the empirical distribution
\begin{equation}
\begin{array}{c|ccccc}
 j_* & 2&3&4&5&6\\
\hline
\text{count} &1&1&3&1&2.
\end{array}
\label{eq:K7_obstruction_index_counts}
\end{equation}
In particular, $j_*=1$ did not occur.  The dominant obstruction in these eight cases is therefore not the first hemispheric condition alone; it occurs in higher symmetric combinations of the seven phases.
\end{experiment}

\begin{remark}
\label{rem:K7_scope}
Experiment~\ref{exp:K7_rigid_obstruction} is a statement about the eight grid points at which the full $K=7$ core was available.  The other $233$ points are not evaluations of $\Delta_7$ and must not be counted as rigidly feasible or infeasible.  Likewise, a reported prefix statistic $\kappa=0$ only excludes the tested angularly ordered prefixes; it does not exclude every possible subfamily of the seven frequencies.
\end{remark}

\begin{corollary}
\label{cor:auxiliary_necessary}
For each configuration in Experiment~\ref{exp:K7_rigid_obstruction} for which a nonnegative Leslie realization containing the same prescribed roots exists in a larger dimension, at least one auxiliary root is necessary.
\end{corollary}

\begin{proof}
Equation~\eqref{eq:K7_all_negative} excludes degree-zero completion, i.e. the rigid characteristic polynomial itself.  Any larger realization containing the same roots therefore has a nonconstant complementary factor.
\end{proof}

\section{Minimal auxiliary completion as an exact linear feasibility problem}
\label{sec:minimal_auxiliary_completion}

The negative rigid test does not determine how many auxiliary eigenvalues are necessary.  That question has an exact finite-dimensional formulation that does not require re-estimating any frequency.

Let
\begin{equation}
C(z)=z^d+c_1z^{d-1}+\cdots+c_d,
\qquad d=2K+1,
\label{eq:core_polynomial_C}
\end{equation}
be a fixed rigid core polynomial $P_{b,\Omega}$, so $C(b)=0$.  For $q\geq0$ let
\begin{equation}
R_q(z)=z^q+r_1z^{q-1}+\cdots+r_q
\label{eq:auxiliary_multiplier}
\end{equation}
be a real monic polynomial and write
\begin{equation}
C(z)R_q(z)
=
z^{d+q}+\sum_{k=1}^{d+q}p_k z^{d+q-k}.
\label{eq:completed_product_coefficients}
\end{equation}

\begin{definition}
\label{def:auxiliary_completion_index}
The auxiliary completion index of the core $C$ is
\begin{equation}
\nu(C)
=
\min\left\{
q\geq0:
\exists R_q\text{ monic real with }p_k\leq0\text{ for every }k\geq1
\right\},
\label{eq:auxiliary_completion_index}
\end{equation}
with $\nu(C)=+\infty$ if no such polynomial exists.  For $C=P_{b,\Omega}$ we also write $\nu_b(\Omega)=\nu(C)$.
\end{definition}

\begin{proposition}
\label{prop:completion_equivalence}
A degree-$q$ auxiliary completion in Definition~\ref{def:auxiliary_completion_index} exists if and only if there is a nonnegative Leslie companion matrix of dimension $d+q$ whose spectrum contains all roots of $C$ and whose Perron eigenvalue is $b$.
\end{proposition}

\begin{proof}
If $p_k\leq0$, set $a_k=-p_k\geq0$.  Then \eqref{eq:completed_product_coefficients} is
\[
z^{d+q}-\sum_{k=1}^{d+q}a_k z^{d+q-k},
\]
the characteristic polynomial of a nonnegative Leslie companion matrix.  Since $C(b)=0$, the product also vanishes at $b$.  The vector
\[
(b^{d+q-1},b^{d+q-2},\ldots,b,1)^T
\]
is strictly positive and is an eigenvector with eigenvalue $b$.  The Collatz--Wielandt characterization therefore gives spectral radius $b$.

Conversely, if a nonnegative Leslie characteristic polynomial contains all roots of $C$, polynomial division by the monic factor $C$ gives a real monic complementary polynomial $R_q$, and nonnegativity of the Leslie coefficients is exactly the condition $p_k\leq0$.
\end{proof}

Set $c_0=1$ and $c_k=0$ outside $0\leq k\leq d$, and set $r_0=1$.  Then
\begin{equation}
p_k
=
\sum_{j=0}^{q}r_jc_{k-j}
=
c_k+\sum_{j=1}^{q}c_{k-j}r_j.
\label{eq:product_convolution}
\end{equation}
Define the $(d+q)\times q$ Toeplitz matrix
\begin{equation}
(A_q)_{kj}=c_{k-j},
\qquad
1\leq k\leq d+q,
\quad 1\leq j\leq q,
\label{eq:completion_toeplitz_matrix}
\end{equation}
and the vector
\begin{equation}
h_q=(-c_1,\ldots,-c_d,0,\ldots,0)^T\in\mathbb R^{d+q}.
\label{eq:completion_rhs}
\end{equation}

\begin{theorem}
\label{thm:completion_lp}
For fixed $q$, a degree-$q$ auxiliary completion exists if and only if the linear system
\begin{equation}
\boxed{
A_q r\leq h_q
}
\label{eq:completion_lp}
\end{equation}
is feasible for $r=(r_1,\ldots,r_q)^T$.  Consequently,
\begin{equation}
\boxed{
\nu(C)=\min\{q\geq0:A_qr\leq h_q\text{ is feasible}\}.
}
\label{eq:nu_first_feasible_lp}
\end{equation}
\end{theorem}

\begin{proof}
Equation~\eqref{eq:product_convolution} gives $p_k\leq0$ exactly as the $k$th inequality in \eqref{eq:completion_lp}.  Definition~\ref{def:auxiliary_completion_index} then gives \eqref{eq:nu_first_feasible_lp}.
\end{proof}

\begin{proposition}[Dual infeasibility certificate]
\label{prop:farkas_completion_certificate}
Degree $q$ is impossible if and only if there exists a vector $y\in\mathbb R^{d+q}$ satisfying
\begin{equation}
\boxed{
y\geq0,
\qquad
A_q^Ty=0,
\qquad
h_q^Ty<0.
}
\label{eq:farkas_completion_certificate}
\end{equation}
Such a $y$ is a Farkas certificate that no choice of $q$ auxiliary roots can repair the rigid core.
\end{proposition}

\begin{proof}
This is Farkas' alternative applied to the system $A_qr\leq h_q$.
\end{proof}

\begin{proposition}
\label{prop:completion_upward_closed}
If degree $q$ is feasible, every degree $q'\geq q$ is feasible.  Hence the feasible auxiliary degrees form either the empty set or the tail
\begin{equation}
\{\nu(C),\nu(C)+1,\nu(C)+2,\ldots\}.
\label{eq:completion_tail}
\end{equation}
\end{proposition}

\begin{proof}
If $R_q$ is feasible, then $zR_q$ is monic of degree $q+1$ and
\[
C(z)zR_q(z)=z\,C(z)R_q(z).
\]
Its nonleading coefficients are the previous nonpositive coefficients shifted by one place, together with a final zero.  Iteration proves the claim.
\end{proof}

\subsection{Auxiliary defect profile}
\label{subsec:auxiliary_defect_profile}

For a quantitative profile before exact feasibility, define
\begin{equation}
E_q(C)
=
\min_{r\in\mathbb R^q,\ e\geq0}
\left\{
e:
A_qr-h_q\leq e\mathbf 1
\right\}.
\label{eq:auxiliary_defect_profile}
\end{equation}
Equivalently, $E_q$ is the minimum possible largest positive coefficient of the completed product.  To compare different $b$ values, the same definition may be applied after the sign-preserving scaling $z=\sqrt b\,w$, which replaces $c_k$ by $c_k/b^{k/2}$.

\begin{proposition}
\label{prop:defect_profile_properties}
The auxiliary defect profile satisfies
\begin{equation}
E_q(C)\geq0,
\qquad
E_{q+1}(C)\leq E_q(C),
\qquad
E_q(C)=0\Longleftrightarrow q\geq\nu(C).
\label{eq:defect_profile_properties}
\end{equation}
Thus $q\mapsto E_q(C)$ records how the rigid obstruction is removed as auxiliary degrees are introduced, while $\nu(C)$ is its first zero.
\end{proposition}

\begin{proof}
Nonnegativity is built into \eqref{eq:auxiliary_defect_profile}.  The construction $R_q\mapsto zR_q$ used in Proposition~\ref{prop:completion_upward_closed} preserves the coefficient violation and proves monotonicity.  Finally, $E_q=0$ is exactly feasibility of \eqref{eq:completion_lp}.
\end{proof}

\subsection{Exact one-auxiliary test}
\label{subsec:one_auxiliary_test}

Because the rigid core has
\begin{equation}
c_d=-b^{K+1}<0,
\label{eq:core_constant_negative}
\end{equation}
a degree-one multiplier $R_1(z)=z+t$ must satisfy $t\geq0$.  Its remaining conditions are
\begin{equation}
c_1+t\leq0,
\qquad
c_k+t c_{k-1}\leq0,
\quad 2\leq k\leq d.
\label{eq:q1_all_conditions}
\end{equation}
Hence $\nu(C)=1$ exactly when the interval of $t$ values determined by \eqref{eq:q1_all_conditions} and $t\geq0$ is nonempty.  Explicitly, every index with $c_{k-1}<0$ supplies a lower bound
\begin{equation}
t\geq -\frac{c_k}{c_{k-1}},
\label{eq:q1_lower_bound}
\end{equation}
every index with $c_{k-1}>0$ supplies an upper bound
\begin{equation}
t\leq -\frac{c_k}{c_{k-1}},
\label{eq:q1_upper_bound}
\end{equation}
and $c_{k-1}=0$ requires $c_k\leq0$.  The first coefficient supplies $t\leq-c_1$.

\begin{remark}
\label{rem:delta_vs_nu}
The rigid margin $\Delta_K$ and the completion index $\nu$ measure different properties.  The sign of $\Delta_K$ decides $q=0$ feasibility, but its magnitude need not determine how many auxiliary degrees are necessary.  The appropriate comparison is therefore not only $-\Delta_K$ versus $\nu$, but the full scaled defect profile
\[
E_0,E_1,\ldots,E_{\nu}=0.
\]
For the eight $K=7$ configurations in Experiment~\ref{exp:K7_rigid_obstruction}, this profile can be computed from the already estimated phase sets without any new frequency estimation.
\end{remark}

\begin{experiment}
\label{exp:minimal_auxiliary_protocol}
For each of the eight evaluable $K=7$ phase sets, form the fixed core polynomial $C=P_{b,\Omega}$ once.  Starting at $q=0$, solve \eqref{eq:completion_lp} in increasing degree.  For every infeasible $q$, retain a dual vector satisfying \eqref{eq:farkas_completion_certificate}; at the first feasible degree retain the multiplier coefficients, the completed Leslie coefficients, the minimum coefficient, the primal residual, and the spectral residual.  Report
\[
(b,\Delta_7,j_*,\nu,E_0,E_1,\ldots,E_\nu)
\]
for every case.  This protocol certifies the minimal auxiliary repair degree rather than merely producing one high-dimensional feasible completion.
\end{experiment}


\section{Decision criterion for the Leslie programme}
\label{sec:decision_criterion}

\begin{conjecture}
\label{conj:transportable_trajectory}
There exists a fixed $K$, a fixed or predeclared base rule, and a low-complexity transport law on $\mathcal M_K\cap\mathbb R_+^m$ such that consecutive rigid Leslie matrices inferred from past Mertens windows satisfy small transport residuals and predict future windows better than the null models without recalibration.
\end{conjecture}

\begin{remark}
\label{rem:falsification}
The conjecture is falsified for a chosen protocol if rigid feasibility occurs only sporadically, mode matching is unstable, finite increments repeatedly leave the admissible set, velocity residuals are no smaller than in null sequences, or the transported model fails to improve future-window prediction.
\end{remark}

\begin{remark}
\label{rem:success}
A positive result would require all of the following: persistent rigid feasibility, coherent branch tracking, compatibility with the corrected cone \eqref{eq:spectral_velocity_cone}, a transport law fixed from past data, and reproducible out-of-sample improvement.  Static feasibility alone is insufficient.
\end{remark}

\section{Conclusions}
\label{sec:conclusions}

\begin{proposition}
\label{prop:final_conclusion}
The moving-window formulation has three exact layers.  First, rigid Leslie spectra possess the strain geometry, positivity constraints, and modal transport developed above.  Second, for fixed $b$ their coefficient vectors lie in the $K$-dimensional reciprocal--Perron affine hull \eqref{eq:reciprocal_perron_hull}.  Consequently overlapping order-$m$ recurrence windows satisfy the exact kernel law \eqref{eq:overlap_kernel_intersection}; the first generic moving stride is $K+2$; and at that stride the update is reduced to one scalar innovation, after which the remaining $K+1$ new recurrence values are causally determined.  Third, when a rigid core is infeasible, its minimal positive spectral repair is the first feasible Toeplitz linear program \eqref{eq:completion_lp}, with exact dual infeasibility certificates \eqref{eq:farkas_completion_certificate}.  These statements are algebraic and do not assume an arithmetic law for Mertens.
\end{proposition}

\begin{remark}
\label{rem:final_status}
The rigid $K=7$ autoconsistency test is already negative on every evaluable grid point in Experiment~\ref{exp:K7_rigid_obstruction}, with a substantial negative margin.  Therefore the rigid moving model cannot be invoked there before repair: auxiliary spectral states are necessary whenever a positive completion exists.  The immediate numerical task is the certified minimal-completion protocol of Experiment~\ref{exp:minimal_auxiliary_protocol}.  Only after identifying the minimal repaired cores does it make sense to test a moving-window transport law.  For any repaired moving model that reaches the critical-stride regime, the remaining irreducible arithmetic quantity is again the scalar innovation $\tau_n$; Leslie geometry alone does not determine it.
\end{remark}

\section*{Notebook History}
Version 1.0: static rigid Leslie family and spectral comparison.\\
Version 2.0: positive spectral completion with auxiliary eigenvalues.\\
Version 3.0: moving rigid Leslie matrices on consecutive windows, corrected tangent-cone formulation, first- and second-order spectral transport, and a falsifiable forecasting protocol.\\
Version 4.0: exact strain coordinates, minimum Perron value, universal critical region, phase Jacobian, viability conditions, and modal amplitude transport.\\
Version 5.0: fixed-Perron reciprocal affine hull, exact overlap rigidity, critical stride $K+2$, scalar innovation transport, admissible innovation intervals, direct determinant tests, and the closed $K=1$ model.\\
Version 6.0: negative rigid $K=7$ margin test, exact minimal auxiliary-completion index, Toeplitz LP formulation, Farkas certificates of lower-degree impossibility, and the monotone auxiliary defect profile.

\begin{thebibliography}{99}

\bibitem{Benvenuti2019}
L.~Benvenuti,
\emph{The inverse eigenvalue problem for Leslie matrices},
Electronic Journal of Linear Algebra \textbf{35} (2019), 319--330.
DOI: \href{https://doi.org/10.13001/1081-3810.3980}{10.13001/1081-3810.3980}.

\bibitem{DuarteKlein2016}
P.~Duarte and S.~Klein,
\emph{Lyapunov Exponents of Linear Cocycles: Continuity via Large Deviations},
Atlantis Studies in Dynamical Systems, vol.~3, Atlantis Press, 2016.
DOI: \href{https://doi.org/10.2991/978-94-6239-124-6}{10.2991/978-94-6239-124-6}.

\bibitem{Yener2023}
N.~Yener,
\emph{On derivatives of eigenvalues, eigenvectors and generalized eigenvectors of matrices},
International Journal of Algebra \textbf{17} (2023), 61--75.
DOI: \href{https://doi.org/10.12988/ija.2023.91740}{10.12988/ija.2023.91740}.

\bibitem{Wright2019}
S.~J. Wright,
\emph{Nonlinear Optimization 2: Lecture Notes},
University of Wisconsin--Madison, 2019.

\end{thebibliography}

\end{document}
